Esta interpretación la redactó Claude a partir de los registros de la corrida; no es del autor.

Astra no llegó nunca a las obligaciones: en los 12 intentos evaluados (episodios 1 y 2 de la corrida original y el episodio 3 repetido) Isabelle rechazó algún lema, y ningún lema quedó verificado, así que la biblioteca de la memoria terminó vacía. La estrategia fue opuesta a la de Luna: en cada intento Astra escribió entre 4 y 11 lemas, casi la demostración completa del teorema partida en piezas (cota del contador, paso del variante, mitad del ancho, signo en el subintervalo, raíz en el interior), en lugar de avanzar lema por lema. Como el juez exige que todos los lemas del bloque se verifiquen antes de mirar las obligaciones, un solo error bastó para perder cada intento. Los primeros errores fueron, sobre todo, demostraciones que no cierran (6) o cuyo método inicial no aplica (3); una vez usó \texttt{field\_simp}, método de otro asistente de prueba, como Luna.

Las llamadas a Astra fueron muy lentas: entre 5 minutos y dos horas y media por intento, parte por tiempos agotados de la API. En el episodio~3 original la cuenta de OpenAI se quedó sin créditos; ese episodio se repitió desde la memoria del final del episodio~2 con el mismo código congelado, y sus 4 intentos sin respuesta del modelo se apartan de todas las métricas.

Frente a Luna con el mismo juez, la misma configuración y el mismo mensaje de memoria, ninguno de los dos aceptó un episodio. Luna verificó 4 lemas y llegó dos veces a las obligaciones; Astra, ninguna. En este objetivo la diferencia entre modelos no está en la matemática del enunciado sino en cuánto arriesga cada intento frente a un juez que no perdona un lema fallido.
